\chapter{Introdução}
\label{chap:introducao}

\orientacao{Uma a duas páginas, nesta ordem: (1) identificar o projeto, o parceiro, o marco contratual e o período; (2) nomear a macroentrega com o texto do Anexo I e dizer o que ela compreende; (3) ligar à entrega anterior, dizendo o que já estava estabelecido e o que esta entrega acrescenta; (4) guiar a leitura, com uma frase por capítulo.

Não abra com contexto de mercado nem com a motivação original do projeto, que pertencem à proposta já aceita. A introdução de uma macroentrega intermediária começa no projeto, não no mundo. Registre em texto que o projeto é realizado com apoio da EMBRAPII, na forma que o contrato exigir.}

O presente relatório documenta a Macroentrega~\ph{NN} do projeto \ph{NOME DO PROJETO}, realizado pela \unidade{} com apoio da EMBRAPII, em parceria com \ph{parceiro tecnológico}, no período de \ph{mês/ano} a \ph{mês/ano}. Esta entrega compreende \ph{descrição da macroentrega, nos termos do Anexo I}.

\ph{Vínculo com a entrega anterior: o que já estava estabelecido na Macroentrega NN-1 e o que esta entrega acrescenta. Obrigatório em macroentrega intermediária.}

O Capítulo~\ref{chap:requisitos} apresenta \ph{o estado dos requisitos e a evolução em relação à entrega anterior}. O Capítulo~\ref{chap:frente1} descreve \ph{o que a frente 1 entregou}. O Capítulo~\ref{chap:discussao} discute \ph{o significado dos resultados e as limitações}. O Capítulo~\ref{chap:conclusao} declara o cumprimento da macroentrega, o nível de maturidade tecnológica atingido e a base deixada para a entrega seguinte.

\section{Objetivos}
\label{sec:objetivos}

\orientacao{Três níveis, sem misturar. O objetivo geral e os objetivos específicos do projeto são transcrição literal do contrato e não mudam de uma entrega para outra: divergência entre relatórios sugere escopo movediço. Os objetivos da macroentrega são o recorte desta entrega dentro deles.

Escreva os objetivos da macroentrega no infinitivo e de modo que o próprio relatório permita verificá-los (``operar o veículo em gramado real'' é verificável; ``avançar na maturidade do sistema'' não é). Se o escopo mudou em relação ao Anexo I, diga aqui, em uma frase, com a referência ao registro da mudança (ata, memorando ou aditivo).

O texto dentro das caixas de transcrição vem do contrato e não se edita, mesmo que a redação pudesse melhorar. Se a fonte não estiver à mão, deixe o marcador e peça o texto à coordenação.}

\subsection{Objetivo Geral do Projeto}

\begin{transcricao}{Anexo I, objetivo geral}
\ph{Transcrever literalmente o objetivo geral do projeto.}
\end{transcricao}

\subsection{Objetivos Específicos do Projeto}

\begin{transcricao}{Anexo I, objetivos específicos}
\begin{enumerate}[label=\alph*)]
    \item \ph{Transcrever literalmente o primeiro objetivo específico};
    \item \ph{Transcrever literalmente o segundo objetivo específico}.
\end{enumerate}
\end{transcricao}

\subsection{Descrição da Macroentrega~\numeromacroentrega{} no Anexo I}

\begin{transcricao}{Anexo I, macroentrega \numeromacroentrega}
\ph{Transcrever literalmente a descrição contratual da macroentrega, igual ao texto do SRInfo.}
\end{transcricao}

\subsection{Objetivos da Macroentrega~\numeromacroentrega}

\ph{Frase de abertura: o que esta macroentrega se propôs a entregar, em redação própria.} Os objetivos são:
\begin{enumerate}[label=\alph*)]
    \item \ph{verbo no infinitivo e objeto verificável pelo próprio relatório};
    \item \ph{verbo no infinitivo e objeto verificável pelo próprio relatório}.
\end{enumerate}
